% Extracted missing sections from Chapter 3
% Footnotes converted to LaTeX format
% OCR artifacts cleaned and normalized

\section[Introduction to the Scholarship]{\origsecnum{3.1}Introduction to the Scholarship}
\label{sec:chapters-ch03:3.1}

IntroductionFor the famous story, first published by Tung in 1930, about Wang Yi-jung's sickness and Liu E's discovery of inscriptions on the ``dragon bones'' from a Peking pharmacy, see Lefeuvre (1975), p. 17, n. 1; Li Chi (1977), pp. 7-11. Liu E himself made no mention of the incident in his introduction to T'ieh-yün. to the Scholarship
Few events have excited the minds of modern Chinese scholars like the discovery
of the oracle bones. In the three-quarters of a century that has elapsed since 1899, when
Wang Yi-jung and Liu E first appreciated the true nature of these ``dragon bone'' inscriptions, a large body of books and articles-at the last count, over 180 books and over 970
monographs has appeared “like bamboo shoots after spring rain.'' In this chapter, I
will first introduce the essential works and indicate how they may be used.Shima (1971), preface, p. 1; cf. Chang Pingch'üan (1967), p. 841, who refers to approximately 1,000 articles and over 180 books by about 300 scholars.³The progress of oracle-bone scholarship (chiaku-hsüeh) has involved several interrelated stages: the identification of the true nature of the inscriptions; the identification of the site from which the inscribed bone fragments were being taken; the collection and publication of those fragments; the eventual conduct of scientific excavations in the Hsiao-t'un region; and, above all, the gradual decipherment, still continuing, of the graphs. The history of this progress has been well summarized elsewhere. See, in particular, Kaizuka (1946), pp. 191-201; (1967), pp. 8-60; Jung Keng (1947); Hu Hou-hsüan (1955), pp. 35-42; Ch'en Meng-chia (1956), pp. 50-54; Tung (1965), pp. 46-118. For Western-language accounts, see Creel (1938), pp. 1-16; Lefeuvre (1975), which is fundamental; and Li Chi (1977).

\section[Reference Works]{\origsecnum{3.3}Reference Works}
\label{sec:chapters-ch03:3.3}

Reference Works
7
OverKaizuka (1967, p. 204, estimates that we know about 3,000 graphs on the oracle bones. about 800 of which can be interpreted; Chang Ping-ch'üan (1967), p. 829, refers to 3,320 graphs, of which we can recognize about 1,000, with 500 to 700 more identifiable but in dispute. Hu Houhsüan (1945), p. 5b, speculated, extravagantly in my view, that the full oracle-bone vocabulary, were we able to recover it, would have included 5,000 to 6,000 words, of which we would be able to recognize from 2,000 to over 3,000. No exactitude is possible in such counts (for other examples, see Chang Ch'un-shu [1972], p. 283); nor is anything more than a general estimate necessary. See too n. 20. three thousand Shang graphs have been identified on the oracle bones, and
about half of these may be translated or identified with varying degrees of certainty. The
dictionaries and concordances listed below are indispensable to the study of the inscriptions.
3.3.1
DictionariesChung-kuo k'e-hsüeh-yüan k'ao-ku yen-chiu-so ‹† 國科學院考古研究所,Chia-ku wen-pien 甲骨文編, K'ao-ku-hsüeh chuan-k'an yi-chung ti-shih-ssu hao 考古學專刋乙種第十四號(Peking,1965). This work, prepared under the general editorship of Sun Hai-po is a revised edition of his Chia-ku wen-pien (Peiping, 1934), which identified some 500 to 600 graphs, taken from eight collections (cf. Teng and Biggerstaff [1971], p. 46). Some of the 4,672 entries in CKWP probably refer to the same words and could be combined. The 1,723 entries in the main section of the dictionary consist of words which have been ``identified'' either in the sense that a Shuo-wen graph can be assigned (in 941 cases) or that a modern graphic reconstruction and putative position within the Shuo-wen system of classification can be provided. A ho-wen section contains 371 entries for graphs which frequently appear in conjunction, such as those used in ancestral titles. The appendix contains 2,949 entries of graphs which have not been deciphered or classified (CKWP, preface, p. 1). The forty collections consulted are listed with their abbreviations between the appendix and the modern-character index. An edition of this dictionary was reissued by Yi-wen in Taipei in 1975 under the title Chiao-cheng chia-ku wen-pien 校正甲骨文編,
8
The standard dictionary is Chia-ku wen-pien (abbreviated ChIP), which contains
a total of 4,672 entries taken from forty collections of published inscriptions. This needs
to be supplemented, on occasion, with Hsü chia-ku wen-pien.''
The entries in both dictionaries, like those in many works of oracle-bone scholarship,
are arranged according to the 540 radicals of Shuo-wen. Each entry consists of the Shuo-wen
graph (when known) as heading, followed by hand-copied reproductions of the oracle-bone characters it is thought to represent. Occasionally, there is brief discussion of the
graph's meaning, but generally the Shuo-wen graph is the only “definition” provided. The

[^9]: Chin Hsiang-heng, Hsü chia-ku wenpien (Taipei, 1959). This is a supplement to Sun Hai-po's 1934 dictionary (referred to in n. 8). Based on thirty-eight published collections, it contains 1,048 oracle-bone graphs identified with Shuo-wen equivalents, 1,585 graphs not found in Shuo-wen, an appendix of ho-wen, and an appendix of some 300 unidentified graphs. There are two modern-character indexes for the graphs with and without Shuo-wen equivalents. Cf. Noel Barnard in RBS 5 (1959), p. 463. All three dictionaries may need to be consulted on occasion. For example, for the graph yüanë, CKWP (1965) has more citations (thirty-two) than the 1934 edition (nine citations) and the Hsü chia-ku wen-pien (twelve citations) combined. But due to the decision of the 1965 editors (preface, p. 3) not to reproduce identical forms of the same character, the 1965 edition has many less entries for the graph chiu than the 1934 edition. 59 11 } U 11 0

\subsection[Dictionaries]{\origsecnum{3.3.1}Dictionaries}
\label{sec:chapters-ch03:3.3.1}

DictionariesChung-kuo k'e-hsüeh-yüan k'ao-ku yen-chiu-so ‹† 國科學院考古研究所,Chia-ku wen-pien 甲骨文編, K'ao-ku-hsüeh chuan-k'an yi-chung ti-shih-ssu hao 考古學專刋乙種第十四號(Peking,1965). This work, prepared under the general editorship of Sun Hai-po is a revised edition of his Chia-ku wen-pien (Peiping, 1934), which identified some 500 to 600 graphs, taken from eight collections (cf. Teng and Biggerstaff [1971], p. 46). Some of the 4,672 entries in CKWP probably refer to the same words and could be combined. The 1,723 entries in the main section of the dictionary consist of words which have been ``identified'' either in the sense that a Shuo-wen graph can be assigned (in 941 cases) or that a modern graphic reconstruction and putative position within the Shuo-wen system of classification can be provided. A ho-wen section contains 371 entries for graphs which frequently appear in conjunction, such as those used in ancestral titles. The appendix contains 2,949 entries of graphs which have not been deciphered or classified (CKWP, preface, p. 1). The forty collections consulted are listed with their abbreviations between the appendix and the modern-character index. An edition of this dictionary was reissued by Yi-wen in Taipei in 1975 under the title Chiao-cheng chia-ku wen-pien 校正甲骨文編,
8
The standard dictionary is Chia-ku wen-pien (abbreviated ChIP), which contains
a total of 4,672 entries taken from forty collections of published inscriptions. This needs
to be supplemented, on occasion, with Hsü chia-ku wen-pien.''
The entries in both dictionaries, like those in many works of oracle-bone scholarship,
are arranged according to the 540 radicals of Shuo-wen. Each entry consists of the Shuo-wen
graph (when known) as heading, followed by hand-copied reproductions of the oracle-bone characters it is thought to represent. Occasionally, there is brief discussion of the
graph's meaning, but generally the Shuo-wen graph is the only “definition” provided. The

[^9]: Chin Hsiang-heng, Hsü chia-ku wenpien (Taipei, 1959). This is a supplement to Sun Hai-po's 1934 dictionary (referred to in n. 8). Based on thirty-eight published collections, it contains 1,048 oracle-bone graphs identified with Shuo-wen equivalents, 1,585 graphs not found in Shuo-wen, an appendix of ho-wen, and an appendix of some 300 unidentified graphs. There are two modern-character indexes for the graphs with and without Shuo-wen equivalents. Cf. Noel Barnard in RBS 5 (1959), p. 463. All three dictionaries may need to be consulted on occasion. For example, for the graph yüanë, CKWP (1965) has more citations (thirty-two) than the 1934 edition (nine citations) and the Hsü chia-ku wen-pien (twelve citations) combined. But due to the decision of the 1965 editors (preface, p. 3) not to reproduce identical forms of the same character, the 1965 edition has many less entries for the graph chiu than the 1934 edition. 59 11 } U 11 0

modern-character indexes in the rear permit the reader to find a graph, provided he
already knows its modern form.
More useful as a dictionary is Chia-ku wen-tzu chi-shih (abbreviated CKWT).Li Hsiao-ting *, Chia-ku wen-tzu chishih X, Chung-yang yen-chiu-yüan lishih yü-yen yen-chiu-so chuan-k'an chih wu-shih 中央研究院歷史語言研究所專利之五十,8 vols. (Nankang, 1965). 10 Based
upon the graphs listed in Chia-ku wen-pien (1934) and Hsü chia-ku wen-pien (1959), and also
organized according to the Shuo-wen system, this handwritten work assembles in one place
the commentaries, sometimes abbreviated, of scholars who have discussed the etymologies
of individual oracle-bone graphs. At the end of each entry, Li Hsiao-ting reviews the
evidence and gives his own conclusions. A table of contents and various indexes to the
main text, a supplement, and first appendix (for graphs not yet identified with certainty)
make the work convenient to consult. The opinions of some scholars have at times been
overlooked, but the work is an essential reference to early Chinese etymology, focusing
on earliest word usage as found in the Shang inscriptions.For more information about CKWT, see my account (forthcoming) in RBS 11. Recent studies of particular graphs may be found in Chung-kuo wen-tzu X‡ (Taipei); there is an index in vol. 41 (1971). This periodical ceased publication, temporarily at least, as of vol. 52 (June 1974). 11 It is not cited specifically on
many occasions in this book because I assume it to be a first source to which any scholar
will turn for information about particular graphs.The second edition of Shima's Sōrui (sec. 3.3.2) keys each entry to the proper page in CKWT and conveniently assembles these cross references at S602-607. 12
Chia-ku-hsüeh cannot be divorced from chin-shih-hsüeh
Shang bronze
inscriptions which contain more than a minimal number of graphs are relatively few,
but since the language of the bronzes may throw light on the language of the bones, and
since our understanding of Shang graphs frequently depends on our understanding of
the graphs found on Chou bronzes, the scholar of Shang history should be prepared to
consult the two major bronze inscription dictionaries: Chin-wen-pien and Chin-wen ku-lin.
Chin-wen-pien (abbreviated CWP) contains in its main text entries for 1,894 graphs found in
Shang and Chou bronze inscriptions, arranged according to the Shuo-wen system.Jung Keng, Chin-wen-pien (Peking, 1959). This is a revised edition of Jung Keng, Chin-wen-pien (Tientsin [?], 1925). The 1959 edition has been reissued by Lien-kuan ch'u-pan she as the first part of Chin-wen chenghsü-pien ho-ting-pen E✰★ (Taipei, 1971). For more information about the work, see Teng and Biggerstaff (1971), p. 47.¹³ Each
entry is headed by the Shuo-wen form or a modern graph reconstruction, followed by the
bronze forms listed in columns. A modern graph equivalent is listed under the first bronze
form with occasional brief discussion. The name of the vessel from which it comes is given
under every graph. An appendix with 1,199 entries gives graphs and insignia whose
modern forms have not been identified. This is followed by an index to the 3,165 vessels
consulted and an index of modern graphs. Like Chia-ku wen-pien, Chin-wen-pien is less a
dictionary than a list of occurrences. Definitions are minimal.
Chin-wen ku-lin, which supersedes Chin-wen-pien, is a combined concordance and

\subsection[Concordances and Compendiums]{\origsecnum{3.3.2}Concordances and Compendiums}
\label{sec:chapters-ch03:3.3.2}

Concordances and Compendiums
Inkyo bokuji sōrui (abbreviated S) is a concordance of the oracle inscriptions in
sixty-three collections. 19 Over three thousand words have been classified. 20 Under each

entry, Shima Kunio has hand copied, in a standardized oracle-bone script (examples of
which grace the margins of this book) all the inscriptions in which the graph in question
appears, giving the citation for each. There are inevitable errors, inconsistent transcriptions, and omissions; some entries for the more common graphs have been abbreviated
(this is noted); and inscriptions published since 1959, together with some minor collections
published only in journals, are not included.It is true that the second edition identifies rejoined Yi-pien fragments that appear in the Ping-pien volumes published down to 1967, i.e., down to Ping-pien 512 (for indexes to these rejoinings, see n. 72); but in terms of new fragments (as opposed to new pieces; cf. appendix 3, n. 1) no collection published since 1959 is included. 21 But the work's fundamental contribution
is threefold. First, it attempts to provide a comprehensive listing of word usage in context.
Second, it provides control over word frequency (except, as already noted, in the case of
common graphs). Third. since Shima copies the inscriptions under each entry in chronological order (cf. sec. 4.1.1), it provides (within the limits imposed by the accuracy of his
transcription) a rapid introduction to the epigraphic evolution, if any, of individual
graphs (cf. sec. 4.3.1.4 and to their periodicity (appendix 5).
The work is also of major importance because, unlike the oracle-bone dictionaries
described above (sec. 3.3.1), it allows a graph entry to be found, even when the seal or modern
form of the graph is not known. This may be done by using the rear index (pp. 590—601),
arranged under the 164 oracle-bone ``radicals'' devised by Shima. Once the entry has
been found, researchers may study all occurrences of the graph in context; they may move
on to the commentaries which have been published for certain of the inscriptions cited;
or they may, in the 840 cases for which the cross reference is provided, move to the relevant
page in Li Hsiao-ting's Chia-ku wen-tzu chi-shih.Examples of how this works are given in secs. 3.6.1 to 3.6.3; see too n. 12. For my fuller description and evaluation of Inkyo bokuji sōrui, see Keightley (1969a) and RBS 12 (forthcoming).22 Inkyo bokuji sōrui and Chia-ku wen-tzu
chi-shih are complementary reference works. Neither should be relied on as a primary
source for the inscriptions or commentaries. But they are invaluable guides to the sources
and scholarship, the two fundamental works in the library of any oracle-bone scholar.
Based on a corpus of fifty-eight published and seven unpublished collections, Yin-tai
chen-pu jen-wu t'ung-k'aoJao Tsung-yi Yin-tai chen-pu jen-wu t'ung-k'ao, 2 vols. (Hongkong, 1959). 23 prints in modern transcription all the divination charges associated
with the names of 137 diviners.Cf. ch. 2, n. 13. 24 Within each diviner's section, the inscriptions are subdivided by topic: rain, the ten-day week, eclipses, hunting, and so on. A compendium
rather than a concordance, Jao Tsung-yi's work is less useful than Shima's; first, because
it is organized by diviner and topic rather than by words; second, because divinations
lacking a diviner's name are not included; third, because his modern-character transcriptions, useful though they are to the beginner, introduce a considerable element of interpretation. But it is a valuable supplement, especially for studying the role of particular
diviners. Once the diviner's name and the topic of a particular charge are known, a modern
transcription may generally be found without difficulty.For a more detailed assessment of Yin-tai chen-pu jen-wu t'ung-k'ao, see Noel Barnard in Monumenta Serica 19 (1960), pp. 480-486. _ 25

\section[Bibliographies]{\origsecnum{3.4}Bibliographies}
\label{sec:chapters-ch03:3.4}

Bibliographies
GENERAL STUDIES
Four major bibliographies of oracle-bone studies have been published. The first
is Hu Hou-hsüan's Wu-shih-nien chia-ku-hsieh lun-chu-mu, arranged topically and with
author and title indexes.Hu Hou-hsüan, Wu-shih-nien chiaku-hsüeh lun-chu-mu + (Shanghai, 1952; reissued, Hongkong, 1966). Hu (1955), pp. 38-41, lists by category the principal books likely to be useful for beginners. 26 Although useful, it has been superseded to a large extent by
Peng's ``An Annotated Bibliography of the Publications on the Inscribed Oracle Bones
Excavated from the Yin Ruins,Madge Shu-chi Peng, ``An Annotated Bibliography of the Publications on the Inscribed Oracle Bones Excavated from the Yin Ruins,'' Chinese Culture 6.3 (1965), pp. 97-149. The annotations are entirely bibliographical in nature.'' 27 arranged by author, and by Tung and Huang's Hsü
chia-ku nien-piao.Tung Tso-pin and Huang Jan-wei , Hsu chia-ku nien-piao (Taipei, 1967). For further discussion of oracle-bone bibliographies, see my account (forthcoming) in RBS 12. 28 This last is arranged chronologically, but the author index in the rear
permits it to be used as a bibliography. The most recent bibliography is Christian Deydier's
Les Jiaguwen: essai bibliographique et synthèse des études; this covers the period down to 1974
and provides romanizations for the works cited.Christian Deydier, Les Jiaguwen: essai bibliographique et synthèse des études, Publications de l'école française d'extrême-orient, vol.

\section[How to Read the Graphs on a Bone or Shel]{\origsecnum{3.6}How to Read the Graphs on a Bone or Shell Fragment}
\label{sec:chapters-ch03:3.6}

How to Read the Graphs on a
Bone or Shell Fragment
Nothing appears more arcane, at first glance, than an oracle-bone rubbing. I
will try to demystify it in this section with a step-by-step analysis of how an actual turtle-shell piece, composed of two fragments (fig. 20), may be identified, deciphered, and placed
in context. The flow chart (fig. 21: summarizes the possible approaches.

\subsection[How to Identify the Fragment]{\origsecnum{3.6.1}How to Identify the Fragment}
\label{sec:chapters-ch03:3.6.1}

t
Orienting the piece is rarely a problem since rubbings are normally printed
correctly.Fragments have been printed upside down, especially in the earliest collections. See, for example, T'ieh-yün 23.4, 28.4, and the other cases listed in Yen Yi-p'ing's preface to the Yi-wen (Taipei, 1959) reprint of T'ieh-yün. Errors of this sort appear with some frequency in secondary works, e.g., Eisenberger (1938), p. 67 (the inscription is probably fake); Ch'en Chih-mai (1966), fig. 4a (whole plastron inverted); William S. Y. Wang (1972), p. 53 (bottom two fragments, left column, inverted). 64 If necessary, orientation is best done either by identifying certain graphs, or
graph elements, whose correct orientation is known, or by orienting the crack numbers
(near the top of the pu crack; sec. 2.4) and crack notations (near the bottom of the pu
crack; sec. 2.6) correctly. In figure 20, the crack number 6 and the shang-chi at the
left are readily placed. The pu cracks all run left; and the crack numbers 6, 7, and 8 run
right; these clues, together with the general configuration of the piece, lead us to suppose
that it comes from the right side of a shell in the region of the hyoplastron or hypoplastron
(fig.See ch. 2, n. 122. 3).65
As we have seen, inscriptions were written from top to bottom (sec. 2.9.4), but it is
necessary to determine in this case, as in all others, whether the inscription should be read
sideways and then down (in fig. 20, graphs 1, 4, 6, 2, 3, 5, 7; or 6, 4, 1, 2, 7, 5

